% ==========================================================================
% GENERATED by scripts/split_paper.py. Do not edit this file.
%
% Source: imported_papers/Koopman-JEPA-TMLR-rewrite/main.tex
% Part 14: Latent-Norm Drift and the Random Orthogonal Control
%
% Edits belong upstream, in the published paper. To change what the
% chapter includes, change imports/bounded.toml and re-run
% the splitter.
% ==========================================================================
\section{Latent-Norm Drift and the Random Orthogonal Control}
\label{app:drift}
% [new + draft] DONE. The random-orthogonal paragraph reports something the
% draft did not: the control beats two of the three learned operators on
% pendulum and loses to the third, which is a more useful result than the draft's
% one-line mention of it.
This appendix checks the claim of Section~\ref{sec:inversion} that the ratio between the two readouts is latent-norm growth, by measuring that growth on the operators instead of inferring it.

\input{tables/table_drift.tex}

Table~\ref{tab:drift} rolls encoded held-out frames forward through each learned operator and reports the median of $\|\K^k\z\| / \|\z\|$. The Cayley variant holds $1$ at every horizon on both environments, which orthogonality requires and which is the arithmetic check that the measurement is doing what it claims. The others separate on pendulum in the order the ratio column of Table~\ref{tab:retrieval} gives, reaching $4.15$ for the soft variant, $691$ for the MLP-JEPA baseline and $6.0\times10^{3}$ for the unconstrained operator at $T{=}20$, and $5.3\times10^{10}$ for the unconstrained operator by $T{=}50$. Cartpole is the weaker case discussed in Section~\ref{sec:inversion}: every operator there stays within a factor of ten at $T{=}20$, and the growth ordering no longer reproduces the ratio ordering.

\paragraph{An operator that satisfies the guarantee without training.}
Substituting a frozen random orthogonal matrix for $\K$, keeping each model's own encoder and probe, gives an operator with exactly the spectral property the Cayley parameterization enforces and no learned dynamics at all. It separates what the constraint is worth from what training under it is worth, and the two answers differ by model. On pendulum at $T{=}20$ the random operator reads $6.1$ in place of the unconstrained operator's $8.6\times10^{5}$ and $6.3$ in place of the soft variant's $20.3$, so for those two, replacing the learned operator with a random orthogonal one improves the probe score. It reads $11.4$ in place of the Cayley variant's $3.68$, so there the training does buy something, by a factor of three. Under the probe, then, most of what the two operators with an unconstrained spectrum had learned was worth less than a random rotation.

\begin{figure}[t]
    \centering
    \includegraphics[width=\linewidth]{paper_figures/fig_compounding_controls.pdf}
    \caption{Error growth against horizon, with the two controls drawn as
    reference lines. \textbf{Left:} pendulum-swingup. \textbf{Right:}
    cartpole-swingup. Probe state MSE, mean over 10 seeds, log scale, with 95\%
    bootstrap intervals as shaded bands. Every series comes from the same
    evaluation pass, so the learned models and the controls share the
    checkpoint, the probe fit and the episodes.}
    \label{fig:compounding_controls}
\end{figure}

\begin{figure}[t]
    \centering
    \includegraphics[width=0.4\linewidth]{paper_figures/fig_eigenvalues.pdf}
    \caption{The learned spectra are consistent with the error growth in
    Table~\ref{tab:drift}. Eigenvalues of each variant's operator on the complex
    plane, all latent dimensions, one representative seed on
    pendulum-swingup. Koopman~(Stable) lies exactly on the unit circle by
    construction, Koopman~(Soft) concentrates near it with a few eigenvalues
    beyond, and Koopman~(Unconstrained) drifts outside it.}
    \label{fig:eigenvalues}
\end{figure}

% =============================================================================
